%==============================================================================
%  CHAPTER 2 — MLOps BACKGROUND AND TECHNOLOGY FOUNDATIONS
%==============================================================================
\chapter{MLOps Background and Technology Foundations}
\label{chap:background}

\section*{Introduction}
This chapter establishes the technical foundations of the project. It defines the
\gls{mlops} discipline and the machine-learning lifecycle, identifies the core
challenges a platform must address, presents the enabling technologies adopted in
this work with the reasoning behind each choice, and finally positions the project
against existing commercial solutions.

% ---------------------------------------------------------------------------
\section{From DevOps to MLOps}
\label{sec:devops-to-mlops}
\gls{devops} unified software development and operations around automation,
continuous integration, and continuous delivery, shrinking the time between
writing code and running it in production. \gls{mlops} extends these principles
to machine-learning systems. The key difference is that a machine-learning system
is not defined by code alone, but by the combination of \textbf{code}, \textbf{data},
and \textbf{model}. Any of the three can change independently and break the system,
which makes classic \gls{cicd} insufficient on its own.

The machine-learning lifecycle the platform automates can be summarized as a loop:

\begin{enumerate}[leftmargin=1.6em]
  \item \textbf{Data}, ingest and prepare a dataset.
  \item \textbf{Train}, run code that produces a model from the data.
  \item \textbf{Track}, record the parameters, metrics, and artifacts of each run.
  \item \textbf{Register}, version the resulting models and manage their stages.
  \item \textbf{Deploy}, serve a chosen model behind an \gls{http} endpoint.
  \item \textbf{Monitor}, observe the running service and its resource usage.
  \item \textbf{Consume}, expose the deployed model through a secured public
        \gls{api} so that external applications, websites, or services can call
        it using an \gls{api} key, without a platform account.
\end{enumerate}

% >>> USER: a simple lifecycle-loop figure helps here.
% \begin{figure}[H]\centering
%   \includegraphics[width=0.8\textwidth]{ml_lifecycle.pdf}
%   \caption{The machine-learning lifecycle automated by the platform}
%   \label{fig:ml-lifecycle}
% \end{figure}

% ---------------------------------------------------------------------------
\section{Core MLOps Challenges}
\label{sec:mlops-challenges}
A platform that automates the lifecycle above must solve several recurring
problems:

\begin{itemize}[leftmargin=1.4em]
  \item \textbf{Reproducibility}, the same code, data, and parameters must
        yield the same run, which requires isolating execution and recording
        inputs.
  \item \textbf{Experiment tracking}, every run's metrics and artifacts must be
        captured automatically and compared over time.
  \item \textbf{Model versioning}, trained models must be named, versioned, and
        moved between lifecycle stages (Staging, Production, Archived).
  \item \textbf{Reliable deployment}, promoting a model to a live service must
        be a one-click, repeatable operation, not a hand-written script.
  \item \textbf{Isolation and scalability}, each training job and each deployed
        model should run in its own isolated, resource-limited unit on shared
        cluster compute.
  \item \textbf{Observability}, the health and resource consumption of the
        cluster and of each running model must be visible.
\end{itemize}

% ---------------------------------------------------------------------------
\section{Enabling Technologies}
\label{sec:technologies}
This section presents the main technologies used and justifies each choice.

\subsection{Containerization and Kubernetes}
\textbf{Containers} package an application with its exact dependencies, solving the
``works on my machine'' problem at the root. \textbf{\gls{k8s}} is the de-facto
orchestrator for containers at scale: it schedules workloads onto nodes, restarts
failed pods, exposes services, and enforces resource limits. The platform runs on
\gls{aks} in production and on KinD (Kubernetes-in-Docker) for local development,
with no code changes between the two. Kubernetes is the natural substrate for
\gls{mlops} because both training jobs and model servers map cleanly onto pods.

\techlogo{kubernetes}{Kubernetes}

\subsection{Pipeline Orchestration: Kubeflow Pipelines and Argo Workflows}
\gls{kfp} lets a training workflow be expressed as a small directed acyclic graph
in Python; it compiles this graph into \textbf{Argo Workflows}, which runs each
step as an isolated pod. This layer is what allows each user-submitted training to
run in \emph{its own} resource-limited pod on cluster compute. Argo handles pod
scheduling, retries, log capture, and artifact passing between steps. In this
project, the pipeline expresses the steps ``download code from object storage
$\rightarrow$ install dependencies $\rightarrow$ run the script $\rightarrow$
autolog results.''

\techlogo{kubeflow}{Kubeflow Pipelines and Argo Workflows}

\subsection{Experiment Tracking and Registry: MLflow}
\textbf{MLflow} is the de-facto open-source experiment-tracking layer for Python.
A single call, \texttt{mlflow.autolog()}, instruments scikit-learn, XGBoost, and
LightGBM to automatically record parameters, metrics, and artifacts with no
further code changes. MLflow also provides a built-in \textbf{Model Registry} with
named versions and Staging / Production / Archived stages, which the platform
surfaces in its own \gls{ui} while keeping MLflow as the source of truth. Being
open source, it carries no vendor lock-in.

\techlogo{mlflow}{MLflow}

\subsection{Model Serving: KServe}
\textbf{KServe} is a Kubernetes operator that turns a declarative
\texttt{InferenceService} resource into a running Deployment, Service, and Ingress,
backed by a framework-specific model server (scikit-learn, XGBoost, PyTorch,
TensorFlow). This reduces ``deploy a model'' to creating a single resource. The
platform uses KServe in \textbf{RawDeployment} mode, which produces a plain
Kubernetes Deployment compatible with a standard NGINX Ingress, rather than the
default Serverless mode, which would require a Knative ingress the cluster does not
provide.

\techlogo{kserve}{KServe}

\subsection{Object Storage: MinIO}
\textbf{MinIO} is an \gls{s3}-compatible object store deployed inside the cluster.
It holds uploaded code and datasets, and all artifacts that MLflow produces.
MinIO is used instead of a managed cloud blob store because it is a true drop-in
for the \gls{s3} \gls{api} that MLflow, KServe, and the AWS SDK all expect,
avoiding the edge-cases of partially-compatible alternatives, while remaining
portable between local and cloud environments with zero code changes.

\techlogo{minio}{MinIO}

\subsection{Database: PostgreSQL}
A single \textbf{PostgreSQL} instance stores both the platform's own tables
(users, projects, pipeline runs, models, deployments) and MLflow's tracking tables.
A single database keeps backup, restore, and operations simple, which is
appropriate for a single-tenant system.

\techlogo{postgresql}{PostgreSQL}

\subsection{Backend: FastAPI}
\textbf{FastAPI} is an asynchronous Python web framework that derives request
validation and an OpenAPI schema directly from type hints. Python is the natural
choice because the entire machine-learning ecosystem, MLflow, the Kubeflow SDK,
the Kubernetes client, is Python, eliminating cross-language marshalling. Its
asynchronous model gives cheap I/O concurrency for database, MLflow, and Kubernetes
calls.

\techlogo{fastapi}{FastAPI}

\subsection{Frontend: Angular}
\textbf{Angular} is a batteries-included framework (router, forms, HTTP, dependency
injection, reactive state) that ships as a \gls{spa}. The project uses Angular~19
with standalone components and signals, the modern idiom, served as static
files by NGINX.

\techlogo{angular}{Angular}

\subsection{Networking: Cloudflare Tunnel}
Because the host environment blocks inbound traffic to the cluster's public
addresses, the platform exposes itself through a \textbf{Cloudflare Tunnel}: a
lightweight agent opens a long-lived \emph{outbound} connection to Cloudflare,
which then forwards inbound \gls{https} traffic to an in-cluster service. This
provides a real \gls{tls} certificate and a public hostname without requiring any
inbound firewall changes.

\techlogo{cloudflare}{Cloudflare Tunnel}

\subsection{CI/CD: GitHub Actions and Azure Container Registry}
Every change pushed to the main branch triggers a \textbf{GitHub Actions} workflow
that builds the backend and frontend container images, pushes them to
\textbf{\gls{acr}}, and upgrades the cluster deployment, delivering continuous
deployment to the public URL.

% ---------------------------------------------------------------------------

\begin{figure}[H]
  \centering
  
\includegraphics[height=1.6cm,keepaspectratio]{github_actions}\hspace{1.5cm}
  \includegraphics[height=1.6cm,keepaspectratio]{acr}
  \caption{Logos of GitHub Actions and Azure Container Registry}
  \label{fig:logo-cicd}
\end{figure}

\section{Comparison with Existing Solutions}
\label{sec:comparison}
Several managed \gls{mlops} platforms exist, notably Amazon
SageMaker~\cite{sagemaker} and Google Cloud Vertex~AI~\cite{vertexai}. They offer
turnkey training, tracking, and deployment, but at the cost of strong vendor
lock-in and recurring fees, and they hide the underlying mechanics. For an
end-of-studies project whose goals include \emph{understanding} the full
\gls{mlops} stack, an open-source, self-hosted assembly is more instructive and
more portable. Table~\ref{tab:comparison} summarizes the trade-off.

\begin{table}[H]\centering
\caption{Managed platforms vs.\ the self-hosted open-source approach}
\label{tab:comparison}
\renewcommand{\arraystretch}{1.3}
\begin{tabularx}{\textwidth}{|>{\raggedright\arraybackslash}p{2.6cm}|X|X|X|}
\hline
\textbf{Criterion} & \textbf{SageMaker / Vertex AI} & \textbf{Databricks} & \textbf{This project (self-hosted)}\\
\hline
Cost            & Pay-per-use, can grow quickly & Subscription + compute & Cluster cost only\\
\hline
Vendor lock-in  & High & High & None (open source)\\
\hline
Control         & Low (managed black box) & Medium & Full\\
\hline
Portability     & Cloud-specific & Cloud-specific & Runs on any Kubernetes\\
\hline
Learning value  & Low (mechanics hidden) & Medium & High (every layer explicit)\\
\hline
\end{tabularx}
\end{table}

\section*{Conclusion}
This chapter defined \gls{mlops}, the lifecycle the platform automates, and the
challenges it must address. It justified each technology in the stack and argued
for a self-hosted, open-source approach over managed alternatives. The next
chapter translates these foundations into concrete requirements and a project
plan.
